%% ED26-REQ-N-d-023-01: complete original source account; proof residuals retained.
\subsection{Lawvere's quantifiers and sheaves: Nice 1970, published 1971}
\label{sec:p10-n-d-023-source-01}

This is the presentation of \emph{Quantifiers and Sheaves}, the six-page
Nice congress article published in the 1971 proceedings
\cite[pp.~329--334]{Lawvere1970}.  Its account of joint work with Myles
Tierney puts geometry first, treats a topology as an internal modality,
and transports constructions for abstract sets to a variable geometric
setting.  Its two applications are Boolean-valued models, including a
continuum-hypothesis argument without transfinite induction in the
displayed construction, and the Chevalley--Hakim global spectrum without
choosing an external site.  These are the article's announced program and
source-specific constructions, not a replacement of set-theoretic
independence by the assertion that every topos is a model of ZFC.

The introductory motivation concerns sets varying with time, space, or a
ring of definition, rather than sets characterized solely by their
elements over the terminal object.  Lawvere interprets opposing
mathematical tendencies through adjoint functors and criticizes reliance
on an externally accumulated power-set hierarchy.  He emphasizes that
such geometric sets usually have intuitionistic logic, need not satisfy
choice, and need not be determined by their global points
\cite[p.~329]{Lawvere1970}.  This motivation does not remove any of the
structural hypotheses in the following definition.

\subsubsection{The original topos package and its internal adjoints}

\begin{hbdefinition}{def:p10-lawvere-sheaves-topos}{Lawvere's original topos package}
In this presentation a topos \(\mathcal E\) has finite limits, finite
colimits, exponentials, and a subobject classifier
\(\mathsf{true}:1\to T\).  For each object \(A\), the functor
\((-)^A\) is right adjoint to \((-)\times A\).  For each monomorphism
\(m:X'\hookrightarrow X\), there is exactly one
\(\chi_m:X\to T\) for which
\[
\begin{array}{ccc}
X'&\longrightarrow&1\\
\downarrow{\scriptstyle m}&&\downarrow{\scriptstyle\mathsf{true}}\\
X&\xrightarrow{\ \chi_m\ }&T
\end{array}
\]
is a pullback.  For \(\varphi:X\to T\), Lawvere writes
\(\{X/\varphi\}\) for the classified subobject.  The source's
truth-value object is \(T\), not the \(\Omega\) used in
Definition~\ref{def:p10-topos}; finite colimits are explicitly part of
this original list of assumptions \cite[pp.~329--330]{Lawvere1970}.
\index{Lawvere, F. William!quantifiers and sheaves}
\index{topos!Lawvere's original package}
\end{hbdefinition}

The article states that \(T\) is an internal Heyting algebra.  Its
completeness is indexed by objects and maps of \(\mathcal E\), not an
assertion that \(\mathcal E\) has every externally indexed limit.  For
\(f:X\to Y\), precomposition induces
\[
 T^f:T^Y\longrightarrow T^X,
 \qquad \exists_f\dashv T^f\dashv\forall_f,
 \qquad \exists_f,\forall_f:T^X\longrightarrow T^Y.
\]
These are order adjunctions between the internal predicate objects.
On classified subobjects the corresponding conditions are
\[
 \exists_f A\leq B\ \Longleftrightarrow\ A\leq f^*B,
 \qquad f^*B\leq A\ \Longleftrightarrow\ B\leq\forall_f A.
\]
For a topological space the sheaf \(T\) has, over an open \(U\),
the open subsets of \(U\), with restriction by intersection.  For
set-valued presheaves on a small \(\mathcal C\), its value at \(c\)
is the set of cribles, or sieves, on \(c\).  Neither example is
generally Boolean.  Formula classification gives actual subobjects for
the appropriate higher-order formulas; it is not a claim that their
truth is determined by a Boolean value at every global point
\cite[p.~330]{Lawvere1970}.

\begin{hbdefinition}{def:p10-lawvere-sheaves-families}{Families, support, and internal presheaves}
An object of \(\mathcal E/X\) is a map \(p:E\to X\), regarded as
an \(X\)-indexed family.  Pullback along \(f:X\to Y\) is
\(f^*:\mathcal E/Y\to\mathcal E/X\).  Its right adjoint, when
written in the source's family notation, is
\[
 \Pi_f:\mathcal E/X\longrightarrow\mathcal E/Y.
\]
The support \(\sigma_X(p)\) is the characteristic predicate of the
image of \(p\); it belongs to \(\mathcal E(X,T)\).

An internal category has an object of objects \(C_0=X\), an object of
arrows \(C_1\), domain, codomain, identity, and composition maps
satisfying the category equations.  A contravariant action on a family
\(E\to X\) gives restriction along each internal arrow, with identity
and composition laws.  Lawvere denotes the resulting category of
internal \(\mathcal E\)-valued presheaves by
\(\widehat C=C^{\mathrm{op}}\mathcal E\).  An internal functor
\(f:C\to C'\) induces restriction
\(f^*:\widehat{C'}\to\widehat C\); the source denotes its
left and right adjoints by \(f_0\) and \(f_*\).
\end{hbdefinition}

\begin{hbtheorem}{thm:p10-lawvere-sheaves-internal-adjoints}{proved=annex:p10.lawvere-sheaves.internal}{Lawvere's announced internal completeness}
For a topos with the original package, the slice pullback functors have
the displayed right adjoints \(\Pi_f\).  The internal presheaf
category \(\widehat C\) is again a topos, and restriction along an
internal functor has both adjoints \(f_0\) and \(f_*\).
The source uses these constructions to obtain the usual exactness
properties and to explain how even an externally countable topos can
have this form of internal completeness \cite[p.~330]{Lawvere1970}.
The original article announces these conclusions without giving their
general constructions or complete proofs.  The annex records that
historical boundary while supplying separate editorial deductions
through the general constructions below.  The concrete
external-small-base case also retains its distinct detailed proofs.
\end{hbtheorem}

\begin{hbdefinition}{def:p10-lawvere-sheaves-functional-graphs}{Power objects and fiberwise functional graphs}
For any object \(W\), the power object is \(T^W\), with
membership the pullback of \(\mathsf{true}\) along evaluation.
Classifier and exponential transposition give, naturally in \(Z\),
\[
 \mathcal E(Z,T^W)\cong\mathcal E(Z\times W,T)
                         \cong\operatorname{Sub}(Z\times W).
\]
These bijections use the original package, not external powersets.
Let \(\delta_W:W\times W\to T\) classify equality.  Let
\(\land:T\times T\to T\) classify the intersection of the
two truth subobjects.  Thus \(a\leq b\) for predicates means
\(a\land b=a\); this equation can be imposed by an equalizer.

For \(p:E\to X\) and \(f:X\to Y\), put
\(D_0=Y\times T^{X\times E}\), with universal membership
predicate \(r(y,R,x,e)\).  Select \(D_1\hookrightarrow D_0\)
by the following equations of predicates, imposed after transposing
in the displayed \(x,e,e'\) variables:
\[
\begin{aligned}
 r(y,R,x,e)\land
   \delta_X(x,p(e))\land\delta_Y(f(x),y)&=r(y,R,x,e),\\
 r(y,R,x,e)\land r(y,R,x,e')\land\delta_E(e,e')
                    &=r(y,R,x,e)\land r(y,R,x,e').
\end{aligned}
\]
They say that the relation lies over the fiber of \(f\), lies
over \(p\), and is single-valued.  They use finite equalizers,
not an assumed quantifier or an assumed unique-choice theorem.
Let \(R_1\hookrightarrow D_1\times X\times E\) be the
restricted membership relation.  Its projection
\(d:R_1\to D_1\times X\) is monic.  Write \(\chi_d\)
for its classifier, and write
\(\gamma(y,R,x)=\delta_Y(f(x),y)\) for the classifier of
\(D_1\times_Y X\hookrightarrow D_1\times X\).
Define \(D\hookrightarrow D_1\) as the equalizer of their
transposes into \(T^X\):
\[
                    D=\operatorname{Eq}
                       (\Lambda\chi_d,\Lambda\gamma).
\]
This final condition makes the functional graph's domain exactly
the fiber, not merely a dense or inhabited part of it.
\end{hbdefinition}

\begin{hbtheorem}{thm:p10-lawvere-sheaves-dependent-product}{proved=annex:p10.lawvere-sheaves.dependent-product}{Dependent products from the original power-object package}
The object \(D\to Y\) above, denoted \(\Pi_f(p)\), has
a canonical evaluation or counit
\[
             \varepsilon_p:D\times_Y X\longrightarrow E,
             \qquad p\varepsilon_p=\operatorname{pr}_X.
\]
For every \(a:Z\to Y\), taking the graph of a map and
evaluating a classified graph give inverse bijections
\[
 \mathcal E/Y(Z,\Pi_fp)
             \cong\mathcal E/X(Z\times_Y X,E),
\]
natural in both the family \(p\) and the test family \(Z\).
Thus \(f^*\dashv\Pi_f\), without an image, exactness, or
choice hypothesis in the construction.

For every pullback square
\[
\begin{array}{ccc}
 X'&\xrightarrow{\ a\ }&X\\
 \downarrow{\scriptstyle f'}&&\downarrow{\scriptstyle f}\\
 Y'&\xrightarrow{\ b\ }&Y,
\end{array}
\]
there is a canonical natural isomorphism
\[
                  b^*\Pi_f(p)\cong\Pi_{f'}(a^*p).
\]
It is the mate whose evaluation is the pulled-back counit, and
it is compatible with identity squares and pasting pullback
squares.  The conclusion concerns actual pullbacks, not arbitrary
commuting squares.  For composable maps,
\(\Pi_{gf}\cong\Pi_g\Pi_f\), and \(\Pi_{1_X}\cong1\).
Every slice \(\mathcal E/X\) satisfies the original topos
package; in particular, its exponential of \(F\to X\) by
\(E\xrightarrow{p}X\) is \(\Pi_p(p^*F)\).
This is an editorial complete construction of the slice assertion
announced on p.~330, not a proof printed in the original article
\cite[p.~330]{Lawvere1970}.
\end{hbtheorem}

\begin{hbtheorem}{thm:p10-lawvere-sheaves-predicate-adjoints}{proved=annex:p10.lawvere-sheaves.predicate-adjoints}{Internal Heyting predicates, stable images, and both quantifiers}
From the original package and the constructed dependent products,
the following statements hold without Boolean negation or choice.
Every map has an epimorphism--monomorphism image factorization,
and these images are stable under arbitrary pullback.  Subobjects
of each \(X\) form a Heyting algebra: intersections are meets,
images of \(A+B\to X\) are joins, and implication is classified
by the truth-value predicate of
\[
       \operatorname{Eq}(\land,\operatorname{pr}_1)
                                        \hookrightarrow T\times T.
\]
Thus \(D\leq(A\Rightarrow B)\) exactly when
\(D\cap A\leq B\), and \(T\) itself is an internal
Heyting algebra.  No complement operation is assumed.

For \(f:X\to Y\) and \(A\hookrightarrow X\),
\(\exists_fA\) is the image of \(A\to Y\), and
\(\forall_fA\) is \(\Pi_f(A\to X)\), regarded as a
subobject of \(Y\).  They satisfy
\[
 \exists_fA\leq B\ \Longleftrightarrow\ A\leq f^*B,
 \qquad f^*B\leq A\ \Longleftrightarrow\ B\leq\forall_fA.
\]
For every pullback square of the preceding theorem,
\[
 b^*\exists_fA=\exists_{f'}a^*A,\qquad
 b^*\forall_fA=\forall_{f'}a^*A,
\]
and \(\exists_f(A\cap f^*B)=\exists_fA\cap B\).

These operations are represented internally by actual maps
\(\exists_f,\forall_f:T^X\to T^Y\).  If
\(r:Z\to T^X\) classifies \(R\hookrightarrow Z\times X\),
their composites with \(r\) classify, respectively,
\[
 \operatorname{Im}(R\to Z\times Y),\qquad
 \Pi_{1_Z\times f}(R\hookrightarrow Z\times X).
\]
The internal order adjunctions therefore hold at every parameter
object \(Z\), not just on the global subobject posets.  Empty
fibers are retained: universal quantification of top is top,
and existential image of the initial subobject is initial.
This proves editorially the indexed completeness of \(T\)
announced on p.~330 \cite[p.~330]{Lawvere1970}.
\end{hbtheorem}

\begin{hbtheorem}{thm:p10-lawvere-sheaves-effective-quotients}{proved=annex:p10.lawvere-sheaves.effective-quotients}{Regular epimorphisms and effective internal equivalence relations}
Every epimorphism in a topos with the original package is the
coequalizer of its kernel pair.  The image factorizations are
therefore regular-epi--mono factorizations stable under pullback.
For any internal equivalence relation
\(R\hookrightarrow X\times X\), transpose its classifier
to the row map \(r:X\to T^X\), and factor it by its
constructed image as
\[
                       X\twoheadrightarrow Q\hookrightarrow T^X.
\]
The first map has kernel pair exactly \(R\) and coequalizes
it universally.  Thus the category is regular and its internal
equivalence relations are effective.  These are explicit exactness
deductions from the constructed package, not exactness assumed in
order to obtain \(\Pi_f\) or its images
\cite[p.~330]{Lawvere1970}.
\end{hbtheorem}

\begin{hbdefinition}{def:p10-lawvere-sheaves-internal-actions}{Internal presheaves with their typed action data}
Let \(C\) be an internal category, with objects \(C_0\), arrows
\(C_1\), source and target \(s,t:C_1\to C_0\), identities
\(i:C_0\to C_1\), and composition \(a b=a\circ b\)
defined when \(t(b)=s(a)\).  All these are maps in
\(\mathcal E\) satisfying the category equations.  An internal
presheaf is a family \(E\to C_0\) with a map
\[
  E\times_{C_0,t}C_1\longrightarrow E,
             \qquad(e,a)\longmapsto e\cdot a,
\]
lying over \(s(a)\), with
\(e\cdot i(c)=e\) and
\((e\cdot a)\cdot b=e\cdot(a b)\).
A morphism is a family map preserving this action.  Write
\(U_C:C^{\mathrm{op}}\mathcal E\to\mathcal E/C_0\)
for the forgetful functor.

The generalized-element notation in the following constructions
denotes actual dependent families: a family of values over the
fiber of \(q\) is \(\Pi_q\) of that family, and a universal
compatibility condition is \(\forall_q\) of the corresponding
equality subobject.  The two preceding proofs construct both.
Thus none of these expressions presupposes an external infinite
product or quantification only over global points.
\end{hbdefinition}

\begin{hbtheorem}{thm:p10-lawvere-sheaves-internal-presheaf-topos}{proved=annex:p10.lawvere-sheaves.internal-presheaf-topos}{The topos of internal presheaves}
For an internal category \(C\), the category
\(C^{\mathrm{op}}\mathcal E\) satisfies the original topos
package.  The functor \(U_C\) creates finite limits and finite
colimits and detects monomorphisms.  It has a free left adjoint
\(L_C\) and a cofree right adjoint \(R_C\), with underlying
families
\[
             U_CL_CB=\Sigma_s(t^*B),\qquad
             U_CR_CB=\Pi_t(s^*B),
\]
where \(\Sigma_s\) means composition of a family map with
\(s\).  On generalized fibers, \(L_CB\) consists of
\((a:d\to c,b\in B_c)\) at \(d\), with precomposition
action; \(R_CB\) at \(c\) consists of all assignments
\((b_a\in B_d)_{a:d\to c}\), also restricted by
precomposition.

For presheaves \(G,F\), their exponential at an internal object
\(c\) consists of the compatible families
\[
 \bigl(\phi_a:G_d\to F_d\bigr)_{a:d\to c},
 \qquad
 \phi_{a b}(g\cdot b)=\phi_a(g)\cdot b
       \quad(b:e\to d).
\]
Evaluation is \((\phi,g)\mapsto\phi_{i(c)}(g)\);
restriction along \(u:c'\to c\) replaces \(\phi_a\)
by \(\phi_{u a}\).
The classifier \(T_C\) is the subfamily of
\(R_C(C_0\times T)\) consisting of assignments
\((v_a\in T)_{a:d\to c}\) with
\[
                         v_a\leq v_{a b}
                              \quad(b:e\to d).
\]
Its truth assignment is constantly true.  A subpresheaf
\(A\hookrightarrow E\) has unique characteristic assignment
\[
                    (\chi_A^C(e))_a=\chi_A(e\cdot a),
\]
where the right-hand classifier is that of its underlying family.
These formulas define internal objects, restrictions, evaluation,
and unique pullback classification, not only fiberwise candidates.
Power objects are \(T_C^G\) with evaluation membership.
This is an editorial proof of the internal-presheaf assertion
announced on p.~330 \cite[p.~330]{Lawvere1970}.
\end{hbtheorem}

\begin{hbtheorem}{thm:p10-lawvere-sheaves-internal-kan}{proved=annex:p10.lawvere-sheaves.internal-kan}{Both adjoints of restriction along an internal functor}
Let \(F:C\to D\) be an internal functor.  Restriction
\(F^*:D^{\mathrm{op}}\mathcal E\to C^{\mathrm{op}}\mathcal E\)
pulls underlying families back along \(F_0\) and acts along
\(F_1\).  For \(A\) on \(C\), its right adjoint has
generalized fibers
\[
 (F_*A)_d=
 \left\{(a_{c,b}\in A_c)_{b:F(c)\to d}:
 a_{c,bF(u)}=a_{c',b}\cdot u
            \text{ for }u:c\to c',\ b:F(c')\to d\right\}.
\]
Dependent products construct the assignments, and equality
subobjects and universal quantification impose compatibility.
Restriction along \(v:d'\to d\) sends the coordinate
\((c,b:F(c)\to d')\) to the old \((c,vb)\) coordinate.

For the left adjoint, form internal \(D\)-presheaves \(B_0,B_1\)
whose generalized elements over \(d\) are respectively
\[
\begin{aligned}
 B_0(d)&:\quad(c,x\in A_c,b:d\to F(c)),\\
 B_1(d)&:\quad(u:c\to c',x'\in A_{c'},b:d\to F(c)).
\end{aligned}
\]
Their action is precomposition of \(b\).  The two maps
\(B_1\rightrightarrows B_0\) send the latter tuple to
\[
              (c,x'\cdot u,b),\qquad(c',x',F(u)b).
\]
Define \(F_0A\) as their coequalizer in the internal presheaf
topos.  This uses its proved finite-colimit construction, not an
assumed externally indexed coend or an assumed effective quotient.
The presheaf argument distinguishes this left-extension notation
from the object component \(F_0:C_0\to D_0\) of the internal
functor.
There are natural inverse bijections
\[
 \operatorname{Hom}_D(F_0A,B)\cong\operatorname{Hom}_C(A,F^*B),
 \qquad
 \operatorname{Hom}_C(F^*B,A)\cong\operatorname{Hom}_D(B,F_*A).
\]
Thus \(F_0\dashv F^*\dashv F_*\), with actual unit/counit
maps and both triangle identities furnished by the inverse
constructions in the annex.  Empty internal index families are
included.  This supplies editorially the internal-functor adjoints
announced on p.~330 \cite[p.~330]{Lawvere1970}.
\end{hbtheorem}

\begin{hbtheorem}{thm:p10-lawvere-sheaves-presheaf-bridge}{proved=annex:p10.lawvere-sheaves.presheaf}{The concrete presheaf instance of the original package}
Fix the small-base and universe hypotheses of
Definition~\ref{def:p10-presheaf-category}.  In
\(\mathcal E=[\mathcal C^{\mathrm{op}},\mathsf{Set}]\), finite limits
and finite colimits are pointwise; \(T=\Omega\), with
\[
 T(c)=\{\text{sieves on }c\},\qquad
 T(u)(S)=\{v:e\to d:u v\in S\}\quad(u:d\to c).
\]
For \(A\subseteq X\), the characteristic map is uniquely
\[
 (\chi_A)_c(x)=\{u:d\to c:X(u)x\in A(d)\}.
\]
The exponential and power object are
\[
 (Y^X)(c)=\operatorname{Nat}(\mathrm y c\times X,Y),\qquad
 T^X(c)\cong\operatorname{Sub}(\mathrm y c\times X).
\]
Exponential restriction is precomposition by
\(\mathrm y u\times1_X\); under the power-object identification
the restriction is pullback by that map.  Evaluation is
\(\operatorname{ev}_c(\alpha,x)=\alpha_c(1_c,x)\), and membership
is \((S,x)\in\in_X(c)\) exactly when \((1_c,x)\in S(c)\).
All restriction, naturality, classification, and uniqueness conditions
hold, including for the empty base.  Thus this category satisfies the
finite-colimit condition in Lawvere's package as well as the host
elementary-topos package.  This is a source-to-host identification of
the concrete example, not attribution of the host's later detailed
proofs to the six-page article.
\end{hbtheorem}

\subsubsection{Why support, dependent product, and choice are different}

\begin{hbdefinition}{def:p10-lawvere-sheaves-choice-conditions}{The two source conditions}
Lawvere separates the following assertions.
\[
\begin{aligned}
(\Pi):\quad&
 \sigma_Y(\Pi_f p)=\forall_f\sigma_X(p)
 &&\text{for every }f:X\to Y\text{ and }p:E\to X;\\
(\exists):\quad&
 E\twoheadrightarrow\{1/\sigma_1(E\to1)\}
 &&\text{has a section for every }E.
\end{aligned}
\]
The first says that taking support commutes with dependent product.
The second splits the epimorphism onto an object's support, not every
epimorphism onto every base.  In the category of formulas of a higher
theory the source calls the latter a Skolem condition.  In geometry an
existential predicate may be true only locally; it need not furnish a
global section \cite[pp.~330--331]{Lawvere1970}.
\index{axiom of choice!support and dependent product}
\index{Skolem condition!Lawvere}
\end{hbdefinition}

\begin{hbtheorem}{thm:p10-lawvere-sheaves-choice}{proved=annex:p10.lawvere-sheaves.choice}{The conjunction, not either condition alone, expresses choice}
In a topos with the slice and predicate adjoints above, the conjunction
of \((\Pi)\) and \((\exists)\) is equivalent to every epimorphism
having a section.  Lawvere calls this geometrically
zero-dimensionality and logically the axiom of choice
\cite[p.~330]{Lawvere1970}.  The annex gives an editorial complete
proof of this equivalence, rather than asserting that the article
prints that proof.
\end{hbtheorem}

\begin{hbexample}{def:p10-lawvere-sheaves-choice-examples}{The source's group and ordered-diagram comparisons}
For permutation representations of a group \(G\), condition
\((\Pi)\) holds: at a base point the dependent product consists of
sections over its fiber, and its support says all the fiberwise
supports hold.  This comparison uses choice in the external category
of sets.  If \(G\) is nontrivial, its regular action has support
\(1\) but no fixed point; its map to \(1\) has no equivariant
section.  Hence \((\exists)\) fails.

For covariant set-valued diagrams on a well-ordered poset \(P\),
\((\exists)\) holds: an inhabited diagram first becomes inhabited
at a least stage, and one element there propagates to all later
stages in its support.  These are the source's \(P\mathsf{Set}\),
equivalently presheaves on \(P^{\mathrm{op}}\), not presheaves
on \(P\) with the opposite restrictions.

The source specifically says \((\Pi)\) fails for the arrow category
\(2\mathsf{Set}\).  Here is an explicit reconstruction.  In the
covariant diagram category \(0\to1\), put
\[
 X=(\{a,b\}\to\{c\}),\qquad
 E=(\{e_a,e_b\}\to\{u,v\}),
\]
where the first map sends both elements to \(c\), the second sends
\(e_a\) to \(u\) and \(e_b\) to \(v\), and the family
\(p:E\to X\) has \(p_0(e_a)=a\), \(p_0(e_b)=b\),
\(p_1(u)=p_1(v)=c\).  It is pointwise surjective.  A section
would have to choose both \(u\) and \(v\) as the image of \(c\),
so none exists.  For \(f:X\to1\), the dependent product has no
element at stage \(0\) but has two at stage \(1\), whereas
\(\forall_f\sigma_X(p)=\mathsf{true}\) at both stages.
The annex verifies these comparisons and their variance
\cite[p.~330]{Lawvere1970}.
\end{hbexample}

\subsubsection{Recursion, geometric morphisms, and associated sheaves}

Lawvere's additional condition \((\omega)\) concerns the forgetful
functor from objects with an endomap to objects:
\[
 \mathcal E^{\Omega}\longrightarrow\mathcal E.
\]
Here the source's \(\Omega\) encodes an endomap; it is not its
truth-value object \(T\).  The functor is required not to be an
equivalence and to have the free-recursion left adjoint
\((-)\times\omega\).  The endomap on the free object is successor
on the \(\omega\) factor.  Truth-value-valued sequences give the
induction case.  The article invokes this condition for coequalizer
analysis and free group or free ring constructions, but intentionally
does not put it in the topos definition.  It announces that it lifts
to a topos defined geometrically over a base satisfying it
\cite[p.~331]{Lawvere1970}.  It supplies no general free-object
construction proof here.

\begin{hbdefinition}{def:p10-lawvere-sheaves-morphisms}{The source's geometric and logical morphisms}
A geometric morphism \(u:\mathcal E'\to\mathcal E''\) is given
by a functor \(u_*:\mathcal E'\to\mathcal E''\) with an exact
left adjoint \(u^*:\mathcal E''\to\mathcal E'\); the left
adjoint preserves finite limits.  A logical morphism preserves, up
to isomorphism, the structures in the original topos package.  The
article uses \emph{local homeomorphism} for a geometric morphism
whose inverse-image part is logical.  This records its stated
usage, without importing a different modern characterization under
the same name \cite[p.~331]{Lawvere1970}.
\index{geometric morphism}\index{logical morphism}
\end{hbdefinition}

\begin{hbtheorem}{thm:p10-lawvere-sheaves-factorization}{proved=annex:p10.lawvere-sheaves.sheaves}{The announced geometric-morphism factorization}
Every geometric morphism
\(u:\mathcal E'\to\mathcal E''\) factors through a topos
\(\mathcal E\) as \(u''u'\), where
\((u'')_*:\mathcal E\to\mathcal E''\) is full and faithful
and \((u')^*:\mathcal E\to\mathcal E'\) reflects
isomorphisms.  The full-and-faithful geometric part is determined
by a topology \(j_u:T''\to T''\), and its domain is the
category of \(j_u\)-sheaves \cite[p.~331]{Lawvere1970}.
The source gives the associated-sheaf route below, not a complete
factorization proof.  The annex supplies a separate editorial
construction using the proved reflector and subobject adjunction
of the given morphism; that does not alter the historical source
standing.  No uniqueness statement stronger than the source's
topology determination is added.
\end{hbtheorem}

\begin{hbdefinition}{def:p10-lawvere-sheaves-topology}{An internal topology and its closure}
The source's topology is a map \(j:T\to T\) satisfying
\[
 j j=j,\qquad j\mathsf{true}=\mathsf{true},\qquad
 j\circ\land=\land\circ(j\times j).
\]
For \(A\hookrightarrow X\) classified by \(\chi_A\), its
\(j\)-closure is classified by \(j\chi_A\).  It is
\(j\)-closed when \(j\chi_A=\chi_A\), and \(j\)-dense
when its closure is \(X\).  An object is separated when its
diagonal is \(j\)-closed.  A \(j\)-sheaf has a unique extension
of each map along a \(j\)-dense monomorphism.  This unpacks the
usual sheaf condition used by the article.

The closure is not Kuratowski closure of subsets of a space.  In
presheaves on the open-set poset, the relevant operator sends an
order ideal of opens to the principal ideal generated by their
union \cite[p.~331]{Lawvere1970}.
\index{Lawvere--Tierney topology}\index{associated sheaf}
\end{hbdefinition}

\begin{hbtheorem}{thm:p10-lawvere-sheaves-closure}{proved=annex:p10.lawvere-sheaves.closure}{The closure properties of the source topology}
The closure defined above is extensive, monotone, idempotent,
preserves finite intersections, and commutes with pullback.
In particular the source's three displayed axioms yield a genuine
universal closure operator on subobjects; an additional arbitrary
external powerset operation is not being assumed.
\end{hbtheorem}

The article's associated-sheaf argument has four explicit observations:
the image \(T_j\) of \(j\) is a sheaf; \(Y^X\) is a sheaf
when \(Y\) is; the closure of \(\Delta_X\) in \(X\times X\)
is an equivalence relation; and, for a monomorphism of \(X\) into
a sheaf \(Y\), the closure of \(X\) in \(Y\) is its associated
sheaf.  Before applying them it uses the singleton map
\(\{\}:X\to T^X\).  It also characterizes sheaves as the
objects that are separated and closed in every separated object
into which they embed.  The reflector is thus obtained without
infinite direct limits, and its exactness is studied through the
maps it inverts.  The annex retains this sequence and explains
the complete editorial deductions below supply those steps; the
source's printed list is not itself a full proof
\cite[p.~331]{Lawvere1970}.

\begin{hbdefinition}{def:p10-lawvere-sheaves-associated-construction}{The closed-singleton construction}
Write \(T_j\hookrightarrow T\) for the equalizer of \(j\)
and \(1_T\); the next theorem identifies it with the source's
image of \(j\).  Put
\[
 R_X=\overline{\Delta_X}\hookrightarrow X\times X,
 \qquad \kappa_X:X\longrightarrow T_j^X,
\]
where \(\kappa_X\) transposes the closed equality predicate
\(j\delta_X\).  It is the \(j\)-closed companion of the
source's singleton map.  Factor it by its already constructed image:
\[
 X\xrightarrow{q_X}S_X\hookrightarrow T_j^X,
 \qquad a_jX=\overline{S_X}^{T_j^X}.
\]
Write \(\eta_X:X\to a_jX\) for the image map followed by
the closure inclusion.  These are finite/internal constructions
using the proved power objects, images and closure, not an
assumed sheaf reflector or external transfinite iteration.
\end{hbdefinition}

\begin{hbtheorem}{thm:p10-lawvere-sheaves-associated-sheaf}{proved=annex:p10.lawvere-sheaves.associated-sheaf}{The four observations and the actual associated sheaf}
The object \(T_j\) is the image of \(j\) and is a
\(j\)-sheaf.  If \(Y\) is a sheaf, \(Y^X\) is a
sheaf for every \(X\).  The relation \(R_X\) is an
equivalence relation and is the kernel pair of \(\kappa_X\)
and \(q_X\); \(S_X\) is separated and is the universal
separated image of \(X\).

The object \(a_jX\) is a sheaf, and precomposition with
\(\eta_X\) gives natural bijections
\[
       \mathcal E(a_jX,Y)\cong\mathcal E(X,Y)
                           \quad(Y\text{ a }j\text{-sheaf}).
\]
Thus the construction defines a functor left adjoint to the
full inclusion of sheaves.  In particular, if
\(X\hookrightarrow Y\) is a mono into any sheaf, its
closure in \(Y\) is an associated sheaf of \(X\).
An object is a sheaf exactly when it is separated and closed
in every separated object into which it embeds.  The assertions
and all their extension/uniqueness arguments hold at arbitrary
generalized parameter objects; no global-point sufficiency or
choice principle is used.  This is a complete editorial proof
of the source's four observations and characterization, whose
printed treatment is a sketch \cite[p.~331]{Lawvere1970}.
\end{hbtheorem}

\begin{hbtheorem}{thm:p10-lawvere-sheaves-sheaf-topos}{proved=annex:p10.lawvere-sheaves.sheaf-topos}{Left exactness, inverted maps, and the sheaf topos}
For \(f:X\to Y\), call the following two conditions
\(j\)-bijectivity: its image is \(j\)-dense in \(Y\),
and \(\Delta_X\hookrightarrow X\times_YX\) is
\(j\)-dense.  Then
\[
 a_j(f)\text{ is an isomorphism}
 \quad\Longleftrightarrow\quad f\text{ is }j\text{-bijective}
 \quad\Longleftrightarrow\quad
 \mathcal E(Y,Z)\to\mathcal E(X,Z)\text{ is bijective
 for every sheaf }Z.
\]
Both conditions are stable under every pullback.  Every unit
\(\eta_X\) is \(j\)-bijective, and the reflector preserves
terminal objects and all pullbacks, hence all finite limits.

The full category \(\mathcal E_j\) of sheaves has finite
limits computed in \(\mathcal E\), finite colimits obtained
by reflecting the base colimits, exponentials computed in
\(\mathcal E\), and subobject classifier
\(\mathsf{true}_j:1\to T_j\).  A mono between sheaves
is an underlying closed mono, with the unique base classifier
factoring through \(T_j\).  Consequently \(\mathcal E_j\)
satisfies the complete original topos package.  The adjunction
\(a_j\dashv i_j\), with full-and-faithful inclusion \(i_j\),
is a geometric morphism whose inverse image is the constructed
left-exact reflector.  No sheafification or sheaf-topos theorem
is assumed as a premise \cite[p.~331]{Lawvere1970}.
For every parameter object \(B\), the same construction in
\(\mathcal E/B\) with topology \(1_B\times j\) has
relative associated sheaf
\[
                    B\times_{a_jB}a_jX\longrightarrow B
                         \quad(X\to B),
\]
using \(\eta_B:B\to a_jB\).  These relative reflectors
commute canonically with every pullback \(B'\to B\),
including noninhabited parameter objects.
\end{hbtheorem}

The geometric example is Godement's construction lifted to an arbitrary
base.  A topology basis has an object of points \(X\), an object of
indices \(A\), and a pairing \(X\times A\to T\).  The
comparison relation on indices is obtained by pulling back inclusion
of subobjects of \(X\) along \(A\to T^X\).  A directness
condition makes the realization/section adjunction
\[
 \mathsf{lim}:A^{\mathrm{op}}\mathcal E
       \rightleftarrows\mathcal E/X:\Gamma
 \qquad (\mathsf{lim}\dashv\Gamma)
\]
a geometric morphism.  Its image is the sheaf category, described
either by a topology on the presheaf side or by a left-exact
cotriple on the family side.  The source does not spell out the
directness predicate or give the complete adjunction construction.
The following reconstruction derives a necessary-and-sufficient
intrinsic criterion for the displayed geometric adjunction, with
separately identified primary comparisons; it is not a claim that
the original article prints the formula
\cite[pp.~331--332]{Lawvere1970}.

\begin{hbdefinition}{def:p10-lawvere-sheaves-godement-basis}{The intrinsic local directness criterion}
Let \(P\hookrightarrow X\times A\) be the pairing's membership
relation, and write \(U_a\hookrightarrow X\) for its indicated
subobject.  The relation \(a\leq b\) means
\(U_a\subseteq U_b\), defined by the internal predicate
\(\forall x\,(P(x,a)\Rightarrow P(x,b))\).
It gives a thin internal category on \(A\).  The source calls
it a poset; if two distinct indices name the same subobject,
the literal pulled-back order is a preorder.  The construction
keeps those indices and their thin category instead of silently
discarding them.

The local directness conditions are the following two epimorphisms:
\[
 P\longrightarrow X,
 \qquad
 \{(x,a,b,c):P(x,c),\ c\leq a,\ c\leq b\}
      \longrightarrow\{(x,a,b):P(x,a)\land P(x,b)\}.
\]
Equivalently, internally every \(x\) lies in some basis member,
and membership in two members has a local common lower member
containing \(x\).  These are existential image/epi conditions,
not choices of globally defined indices.  They remain valid under
every parameter pullback.  For a thin category the parallel-arrow
part of filtering is automatic.  This is the basis special case
of the explicitly epimorphic-local filtering conditions of
\cite[VII\S8, Definition~1 and Lemma~3, pp.~394--397]
{MacLaneMoerdijk1992}; the theorem below proves necessity as well
as sufficiency here, without assuming that book's externally
cocomplete setting.
\end{hbdefinition}

\begin{hbtheorem}{thm:p10-lawvere-sheaves-godement}{proved=annex:p10.lawvere-sheaves.godement}{The intrinsic Godement realization, image, and cotriple}
For the pairing and thin category above, put
\[
       (\Gamma M)_a=\Pi_{U_a\to1}(M|_{U_a})
                  \quad(M\to X),
\]
interpreted internally as \(\Pi_{P\to A}(P\times_XM)\)
over \(A\), with restriction by precomposition on the basis
inclusions.  For a presheaf \(F\) on \(A\), form in
\(\mathcal E/X\) the coequalizer
\[
 H_1(F)\rightrightarrows H_0(F)\longrightarrow L(F),
\]
where generalized elements are
\[
 H_0(F): (x,a,v\in F_a),\ P(x,a),
 \qquad H_1(F):(x,a\leq b,v\in F_b),\ P(x,a),
\]
and the parallel maps give \((x,b,v)\) and
\((x,a,v|a)\).  Then \(L\dashv\Gamma\), with actual
section/germ unit and evaluation counit.  It is the source's
\(\mathsf{lim}\dashv\Gamma\) adjunction.  The two local
directness conditions are equivalent to \(L\) being left exact.
Under them, two germs agree exactly when their representatives
agree on a local common lower basis member.

Its geometric image is the sheaf topos for the internal coverage
in which a sieve \(S\) on \(a\) covers precisely when
\[
                           U_a=\bigvee_{b\in S}U_b.
\]
Here the join is an internally constructed image, not an assumed
external infinite union.  Equivalently the image is the coalgebra
category of the actual left-exact cotriple \(L\Gamma\) on
\(\mathcal E/X\), with counit evaluation and comultiplication
\(L\) of the adjunction unit.  The finite equalizer proof of
this comparison is supplied in the annex, not assumed as a
comonadicity theorem.
Here a coalgebra is a family with \(\theta:M\to L\Gamma M\)
satisfying \(\varepsilon_M\theta=1_M\) and
\((L\Gamma\theta)\theta=(L\zeta_{\Gamma M})\theta\),
where \(\zeta\) is the adjunction unit; its morphisms commute
with these structure maps.

For ordinary sets and a genuine topological basis, these local
conditions are the usual cover/common-refinement conditions and
the coequalizer is the germ construction.  The primary comparison
is Godement's \emph{espace \'etal\'e} construction and generated
sheaf \cite[II\S1.2, pp.~110--112]{Godement1998Printing}.
The held book is a 1998 printing with 1958 copyright; no first-edition
byte identity or printed 1970 directness formula is inferred.
This reconstruction lifts the source's named construction to the
arbitrary base and retains its original announcement standing
\cite[pp.~331--332]{Lawvere1970}.
\end{hbtheorem}

\begin{hbexample}{def:p10-lawvere-sheaves-godement-failures}{Both directness clauses matter}
For \(X\) a singleton and \(A\) empty, realization sends the
terminal presheaf to the empty family, not to \(X\); it does
not preserve the terminal object.  For
\(X=\{x,y,z\}\), \(U_a=\{x,y\}\),
\(U_b=\{x,z\}\), with only these two incomparable indices,
the basis covers \(X\) but has no member refining their
intersection at \(x\).  Realization of the two representables
is \(U_a,U_b\), whose product over \(X\) is \(\{x\}\),
whereas their presheaf product and its realization are empty.
Thus cover alone does not give left exactness.  In a non-set base,
the repaired condition requires epimorphic local witnesses, not
a choice of one index over each global point.
\end{hbexample}

\subsubsection{Double negation and generalized ultraproducts}

Double negation \(j=\neg\neg\) is a topology in any topos.  The
source interprets it as cofinal locality.  Its sheaves have classical
logic: their classifier satisfies \(1+1\cong T\).  This is
equivalent to \(T\) being an internal Boolean algebra and to
each mono having a unique complementary summand
\(X'+(\neg X')\cong X\).  These are source assertions about
the double-negation sheaf topos, not assertions that the original
\(\mathcal E\) was already classical
\cite[p.~332]{Lawvere1970}.

\begin{hbtheorem}{thm:p10-lawvere-sheaves-double-negation}{proved=annex:p10.lawvere-sheaves.double-negation}{Double-negation sheaves have classical internal logic}
The operation \(j(p)=\neg\neg p\) on the internal Heyting
classifier is idempotent, preserves truth and binary meets, and
is therefore a topology.  Its fixed-predicate object \(T_j\)
is the classifier of the constructed sheaf topos.  For closed
subobjects \(A,B\) of a sheaf \(X\), their meet is ambient
intersection, their join is \(\overline{A\lor B}\), and
their complement is the ambient Heyting negation \(\neg A\),
which is itself closed.  Thus every internal predicate in
\(\mathcal E_{\neg\neg}\) satisfies excluded middle,
at every generalized parameter base.

In a topos with the original package, the following conditions
are equivalent: the classifier is an internal Boolean algebra;
every mono has a unique complementary summand; and the canonical
truth/false map \(1+1\to T\) is an isomorphism.  Applying
this to \(\mathcal E_{\neg\neg}\) gives the exact classicality
condition \(1+1\cong T_j\) announced by Lawvere
\cite[p.~332]{Lawvere1970}.  Its coproduct is the one in the
sheaf topos, not an assumed unchanged ambient coproduct.
Classicality alone does not supply the source's additional
Skolem condition in its stronger definition of a Boolean topos,
nor an axiom of choice.
\end{hbtheorem}

\begin{hbdefinition}{def:p10-lawvere-sheaves-ultraproduct}{The source's generalized ultraproduct data}
Take a finite-limit-preserving functor
\(u_*:\mathcal E\to\mathcal E_0\), not necessarily a
geometric morphism, and a homomorphism of internal Heyting algebras
\(h:u_*(T)\to T_0\).  Form \(\mathcal E_h\) by formally
inverting the monomorphisms \(m:X'\hookrightarrow X\) for which
\[
 h\bigl(u_*(\sigma_1(\Pi_{X\to1}(m)))\bigr)
                  =\mathsf{true}_0.
\]
The expression inside \(u_*\) is a predicate \(1\to T\):
it is support of the family of sections of \(m\), or equivalently
universal quantification of its membership predicate.  Finite-limit
preservation identifies \(u_*(1)\) with \(1_0\), making the
display typed \cite[p.~332]{Lawvere1970}.
\index{ultraproduct!Lawvere's generalized construction}
\end{hbdefinition}

\begin{hbtheorem}{thm:p10-lawvere-sheaves-ultraproduct}{sketched=annex:p10.lawvere-sheaves.ultraproduct}{The generalized-ultraproduct announcement}
The localized category \(\mathcal E_h\) is a topos, and
\(\mathcal E\to\mathcal E_h\) is a logical morphism.  Its
relation to \(\mathcal E_0\) is one of closed categories, not
in general a geometric morphism of topoi.  Lawvere invokes this
construction to collapse a Boolean-valued model of abstract sets
to a two-valued model, while doing the principal independence work
before choosing \(h\) \cite[p.~332]{Lawvere1970}.
This is an announced theorem without a localization proof in the
article; no geometric direct-image structure is inferred merely
from preservation of finite limits.
The full ordinary localization and logical-preservation proofs are
\cref{thm:p10-lawvere-localization-denominators,
thm:p10-lawvere-localization-fractions,thm:p10-lawvere-localization-logical}.
The actual full-datum closed functor is constructed in
\cref{thm:p10-lawvere-closed-base-partials,thm:p10-lawvere-closed-base-closed};
it is not automatically strong or normal, and its enriched global
underlying category is not silently identified with this ordinary
localization. The nondegenerate Boolean-valued collapse, including
existence of the required \(h\) and preservation of choice, is
\cref{thm:p10-lawvere-closed-base-collapse}. A degenerate Boolean
model cannot collapse to nondegenerate \(2\). The original combined
announcement retains its historical sketch boundary and unprinted
closed-base convention, not a claim that these separate editorial
constructions remain missing.
\end{hbtheorem}

For comparison, if \(\mathcal E=\mathsf{Set}^I\), take
\(u_*\) to be product and let a given ultrafilter \(U\) define
\(h:2^I\to2\).  For \(A_i\subseteq X_i\), the displayed
condition becomes
\[
                     \{i:A_i=X_i\}\in U.
\]
This calculation explains the ultrafilter terminology.  It does not
identify a localized category of families with the single ordinary
model-theoretic structure \(\prod_i X_i/U\), nor supply an
ultrafilter without the relevant external existence assumption.

%% New owned fragment; released source01 manuscript files remain frozen.
\subsubsection{The exact selected-mono localization and its logical functor}
\label{sec:p10-lawvere-localization}

The original article's datum is the finite-limit-preserving
\(u_*:\mathcal E\to\mathcal E_0\) and the internal Heyting
homomorphism \(h:u_*(T)\to T_0\), not an assumption of a
geometric inverse image or an ordinary model-theoretic ultrafilter
\cite[p.~332]{Lawvere1970}.  The following construction supplies
the ordinary localization and its logical preservation at exactly
that datum.  Its reduction to a truth filter is proved by the
universal property; no finite, principal, Boolean or prime-filter
restriction is imposed.  The base-relative closed-category and
collapse assertions remain distinct below.

\begin{hbdefinition}{def:p10-lawvere-localization-filter}{The full truth filter induced by the original datum}
Identify \(\operatorname{Sub}_{\mathcal E}(1)\) with the
maps \(1\to T\) by classification.  Define
\[
 \theta_h(U)=h\circ u_*(\chi_U):1_0\longrightarrow T_0,
 \qquad
 \mathcal F_h=\{U\hookrightarrow1:\theta_h(U)=\mathsf{true}_0\},
\]
using the finite-limit identification \(u_*(1)\cong1_0\).
The value \(\theta_h(U)\) is not asserted to be the
characteristic map of \(u_*(U)\).  For a mono
\(m:A\hookrightarrow X\), put
\[
         \gamma_m=\forall_{X\to1}(m)\in\operatorname{Sub}(1),
 \qquad \mathcal W_h=\{m:\gamma_m\in\mathcal F_h\}.
\]
The predicate of \(\gamma_m\) is exactly
\(\sigma_1(\Pi_{X\to1}(m))\), by the proved restriction
of the family adjunction to monos.  Thus \(\mathcal W_h\)
is precisely the monomorphism class in the source display,
without interchanging \(u_*\) with universal quantification.
\end{hbdefinition}

\begin{hblemma}{thm:p10-lawvere-localization-denominators}{proved=annex:p10.lawvere-localization.denominators}{The source denominators have uniform cofinal restrictions}
The map \(\theta_h\) is a Heyting homomorphism on global
truth algebras.  The set \(\mathcal F_h\) contains top,
is upward closed, and is closed under binary intersection.
It is proper when \(\mathcal E_0\) is nondegenerate;
the degenerate-base case permits the improper filter.
For every \(U\in\mathcal F_h\), the mono
\(i_{U,X}:U\times X\hookrightarrow X\) lies in
\(\mathcal W_h\).  Conversely every \(m\in\mathcal W_h\)
contains \(\gamma_m\times X\), and its pullback to
that admitted subterminal is an isomorphism.
Hence localization at the source's full mono class and
localization at all \(i_{U,X}\), \(U\in\mathcal F_h\),
have the same universal property.  This is a mathematical
equivalence of the exact classes, not replacement of the
general datum by a special finite filter.
\end{hblemma}

\begin{hbdefinition}{def:p10-lawvere-localization-fractions}{Fractions and their actual equality relation}
Fix the locally small/universe convention, so \(\operatorname{Sub}(1)\)
and each indicated hom collection are sets.  The objects of
\(\mathcal E/\mathcal F_h\) are those of \(\mathcal E\).
A morphism \(X\to Y\) is represented by
\((U,f)\), where \(U\in\mathcal F_h\) and
\(f:U\times X\to Y\).  Set
\[
 (U,f)\sim(V,g)\quad\Longleftrightarrow\quad
 \text{some }W\in\mathcal F_h, W\leq U\cap V,
       \text{ has }f|_{W\times X}=g|_{W\times X}.
\]
Write \([U,f]\) for its class.  Identity has denominator
\(1\).  Compose \([U,f]:X\to Y\) and
\([V,g]:Y\to Z\) on \(W=U\cap V\): restrict
\(f\) to \(W\times X\), pair it with the unique
parameter map to \(V\), then apply \(g\).
No representatives outside the admitted subterminal are
chosen or required.  The canonical functor is
\(Q(X)=X\), \(Q(f)=[1,f]\).

In the 1992 primary comparison, a representative is a
map \(X\times U\to Y\times U\) over \(U\).
Its first coordinate is the present representative;
the other coordinate is forced because \(U\) is
subterminal.  Product symmetry gives the exact typed
bridge.  The ordinary set-sized quotient of hom sets
does not assume an infinite colimit inside \(\mathcal E\)
\cite[V\S9, pp.~258--260]{MacLaneMoerdijk1992}.
That source's filter convention excludes zero; the
improper-filter boundary here is explicitly separate.
\end{hbdefinition}

\begin{hbtheorem}{thm:p10-lawvere-localization-fractions}{proved=annex:p10.lawvere-localization.fractions}{The complete fraction calculus and localization property}
The displayed relation is an equivalence relation and
the composition is independent of representatives and
refinements, associative and unital.  The functor \(Q\)
inverts every source mono in \(\mathcal W_h\).
For any category \(\mathcal D\) and functor
\(H:\mathcal E\to\mathcal D\) inverting that class,
there is a unique extension on the same objects,
\[
 \bar H[U,f]=H(f)\,H(i_{U,X})^{-1}.
\]
Natural transformations extend uniquely by their original
components.  Consequently the constructed category is
the ordinary \(\mathcal E_h\) formally localized at
the exact class stated on p.~332.  Every source fraction
has a uniform representative of this form, and equality
is exactly agreement on a further admitted common
restriction; there is no additional unproved fraction
existence or coherence requirement.
\end{hbtheorem}

\begin{hbtheorem}{thm:p10-lawvere-localization-logical}{proved=annex:p10.lawvere-localization.logical}{The ordinary localized topos and full logical preservation}
The constructed \(\mathcal E_h\) has finite limits,
finite colimits, exponentials and subobject classifier
\(Q(\mathsf{true}):Q1\to QT\), and \(Q\) preserves
all these structures up to their canonical isomorphisms.
In particular \(Q(Y^X)\) is the exponential of
\(QY\) by \(QX\), with its actual evaluation, and
every localized mono has a unique characteristic map,
not merely those arising from initially chosen strict
monos.  Every finite diagram of germs is handled on
a common admitted restriction, including equalizers
and coequalizers of genuinely localized arrows.

Power objects, membership, the graph constructions of
dependent products, and both internal quantifier maps
are preserved as well.  Thus \(Q\) is a logical
morphism in the original source's package.  The
improper filter gives the degenerate topos and is
included explicitly.  The later primary comparison is
\cite[V\S9, Theorem~2 and Lemma~3, pp.~261--263]
{MacLaneMoerdijk1992}; its proof boundaries and proper
filter convention are not silently promoted to a
different 1970 source statement.
\end{hbtheorem}

\begin{hbexample}{def:p10-lawvere-localization-base-functor}{The original base functor need not descend unchanged}
Let \(\mathcal E=\mathsf{Set}\times\mathsf{Set}\),
\(\mathcal E_0=\mathsf{Set}\), \(u_*\) be binary
product, and \(h:2\times2\to2\) project to the first
coordinate.  These satisfy the original hypotheses.
The subterminal \(U=(1,0)\) is admitted, so its mono
into \(1=(1,1)\) is inverted.  But its image under
\(u_*\) is \(0\to1\), not an isomorphism.
Therefore \(u_*\) itself cannot factor through this
localization.  A new base-relative construction is
needed; finite-limit preservation alone does not
produce a geometric direct image or a strong
closed-functor descent.
\end{hbexample}

\begin{proofboundary}{def:p10-lawvere-localization-boundary}{The exact completed and still-owed source assertions}
The ordinary localization and its logical topos preservation are
proved at the full original \(u_* / h\) datum, through the exact
cofinal reduction to its arbitrary induced filter. The article
announces these conclusions. The additional full-datum closed
functor is constructed in
\cref{thm:p10-lawvere-closed-base-partials,thm:p10-lawvere-closed-base-closed},
without asserting unchanged descent of \(u_*\), an invertible
closed comparison, or arbitrary normality. BVM classification and
choice at the declared cell regimes are
\cref{thm:p10-lawvere-bvm-presentation,thm:p10-lawvere-bvm-realization,thm:p10-lawvere-bvm-correspondence};
the nondegenerate collapse, with its actual \(h\)-existence and
choice hypotheses, is \cref{thm:p10-lawvere-closed-base-collapse}.
The degenerate/improper boundary is not a nondegenerate two-valued
collapse. General stronger normal-relative identification is not
entailed by the original datum; the source's unprinted convention
is kept distinct. CH/ccc/no-epi/transfer, spectrum and independent
recursion/free-algebra proofs remain owed. No finite, principal
or improper special case discharges those separate obligations.
\end{proofboundary}


\subsubsection{Abstract sets, Boolean-valued models, and the CH sketch}

\begin{hbdefinition}{def:p10-lawvere-sheaves-bvm}{The source's abstract-set and Boolean packages}
Lawvere takes the category \(\mathcal S\) of abstract sets to
satisfy \((\Pi)\), \((\exists)\), \((\omega)\), classical
logic \((\neg):1+1\cong T\), and irreducibility
\[
 (\lor):\quad\varphi_1,\varphi_2:1\to T,
 \quad\varphi_1\lor\varphi_2=\mathsf{true}
 \quad\Longrightarrow\quad
 \varphi_1=\mathsf{true}\text{ or }\varphi_2=\mathsf{true}.
\]
Classicality plus irreducibility implies that \(1\) has at most the
two truth/false subobjects; they are distinct when the topos is
nondegenerate.  Exact two-valuedness includes that nondegeneracy
condition.  The converse is not asserted: the source uses
\(M^{\mathrm{op}}\mathsf{Set}\) for a monoid that is not a
group to distinguish it.

In this article a \emph{Boolean topos} means a topos satisfying
both \((\exists)\) and \((\neg)\).  Its classifier is usually
written \(2=T\).  A Boolean-valued model
\(\mathcal E\in\mathsf{BVM}/\mathcal S\) is such a Boolean
topos with a specified geometric morphism to \(\mathcal S\).
This is stronger terminology than merely saying that a topos's
internal logic is Boolean \cite[pp.~332--333]{Lawvere1970}.
\index{Boolean-valued model!Lawvere's definition}
\end{hbdefinition}

\begin{hbtheorem}{thm:p10-lawvere-sheaves-generators}{proved=annex:p10.lawvere-sheaves.generators}{Subterminal generators under the source's Boolean package}
If every subobject has a complementary summand and
\((\exists)\) holds, the subobjects of \(1\) form a generating
family: unequal maps \(f,g:E\to F\) are distinguished by
some map \(U\to E\) with \(U\hookrightarrow1\).
This is the generation conclusion stated on p.~333, with an
editorial full proof in the annex.
\end{hbtheorem}

The article further announces that every Boolean-valued model over
\(\mathcal S\) satisfies \((\Pi)\), hence choice, and that
the bicategory \(\mathsf{BVM}/\mathcal S\) is equivalent to
\(\mathsf{CBA}(\mathcal S)\), the category of
\(\mathcal S\)-complete Boolean algebra objects
\cite[p.~333]{Lawvere1970}. These conclusions are announced, not
proved in the short article. The following editorial construction
supplies the actual functors, comparisons and choice proof while
declaring its morphism/2-cell directions explicitly; their names
alone would not supply those data.

% Adopted Lawvere continuation; working spec fragment preserved separately.
\subsubsection{The abstract-set base and the Boolean-valued correspondence}

Lawvere's Nice 1970 article, published in 1971, announces on p.~333
that its Boolean-valued models satisfy choice and correspond to complete
Boolean algebras in the abstract-set base \cite[p.~333]{Lawvere1970}.
Here ``Boolean'' retains the source's conjunction of classical logic and
support splitting, not just the modern classical-logic condition.
The following constructions are independent editorial proofs of that
announcement at an explicitly specified morphism scope. They do not turn
the six-page source into a source-native detailed proof.

\begin{hblemma}{thm:p10-lawvere-bvm-base}{proved=annex:p10.lawvere-bvm.base}{Bounded well-ordering in the source's base}
Let \(\mathcal S\) have the abstract-set package of
Definition~\ref{def:p10-lawvere-sheaves-bvm}. Then \(\mathcal S\)
is nondegenerate, \(1\) is a generator, every epimorphism splits,
and every object admits an internal well-order. The latter construction
uses power objects and a choice function on nonempty subobjects of that
object; it does not require a cumulative universe, Replacement, an
external ordinal, or a comparison with all models of ZFC.
\end{hblemma}

\begin{hbdefinition}{def:p10-lawvere-bvm-matrices}{Complete Boolean algebras and their matrices}
An \(\mathcal S\)-complete Boolean algebra \(B\) is a Boolean
algebra object with \(\bigvee:\mathcal P_{\mathcal S}B\to B\)
satisfying the least-upper-bound property internally. Meets are obtained
by complementation. A morphism preserves Boolean operations and all
\(\mathcal S\)-indexed joins.

The category \(\operatorname{Mat}(B)\) has objects \((I,b,e)\),
where \(I\in\mathcal S\), \(b:I\to B\), and
\(e:I\times I\to B\) satisfy
\[
 e_{ii}=b_i,\quad e_{ij}\le b_i\wedge b_j,\quad
 e_{ij}=e_{ji},\quad e_{ij}\wedge e_{jk}\le e_{ik}.
\]
A map \((I,b,e)\to(J,c,f)\) is a matrix \(r:I\times J\to B\)
with
\[
\begin{gathered}
 r_{ij}\le b_i\wedge c_j,\qquad \bigvee_j r_{ij}=b_i,\qquad
 r_{ij}\wedge r_{ik}\le f_{jk},\\
 e_{ii'}\wedge r_{ij}\le r_{i'j},\qquad
 r_{ij}\wedge f_{jj'}\le r_{ij'}.
\end{gathered}
\]
These are saturated, total, single-valued relations, not functions between
chosen representatives. Identity is \(e\); composition is
\((sr)_{ik}=\bigvee_j r_{ij}\wedge s_{jk}\).
All displayed indexing and inequalities are internal to \(\mathcal S\).
\end{hbdefinition}

\begin{hbtheorem}{thm:p10-lawvere-bvm-presentation}{proved=annex:p10.lawvere-bvm.presentation}{The actual Boolean-algebra presentation of a model}
For a source Boolean-valued model \(p:\mathcal E\to\mathcal S\),
put \(B_{\mathcal E}=p_*T_{\mathcal E}\). This is an
\(\mathcal S\)-complete Boolean algebra, and there is an equivalence
over \(\mathcal S\)
\[
             \mathcal E\simeq\operatorname{Mat}(B_{\mathcal E}).
\]
Each object is represented by an actual effective quotient of a subobject
of \(p^*I\), for an object \(I\) of the base. The functors on maps,
inverse comparison maps, and their triangle identities are supplied in
the annex. No global points of \(\mathcal E\) are assumed sufficient.
\end{hbtheorem}

\begin{hbtheorem}{thm:p10-lawvere-bvm-realization}{proved=annex:p10.lawvere-bvm.realization}{Realization and choice}
For every \(\mathcal S\)-complete Boolean algebra \(B\), let
\(\operatorname{Sh}_{\mathcal S}(B)\) be internal sheaves on the
poset \(B\), where a sieve on \(b\) covers precisely when its join
is \(b\). Its geometric base map \(q\) has
\(q^*=a\Delta\), and
\[
 q_*T\cong B,\qquad
 \operatorname{Sh}_{\mathcal S}(B)\simeq\operatorname{Mat}(B).
\]
It is classical and every epimorphism splits. Therefore it is a Boolean
topos in the source's stronger sense and satisfies \((\Pi)\) and choice.
The same splitting construction applies to every \(\mathcal E\) in
the preceding theorem. It uses a well-order of a base indexing object
and Boolean disjointification, not a choice of a sheaf-global witness.
The one-element Boolean algebra is allowed and realizes the degenerate
topos; excluding it on one side requires excluding it on the other.
\end{hbtheorem}

\begin{hbdefinition}{def:p10-lawvere-bvm-cell-regime}{The specified morphism and 2-cell regimes}
Write \(\mathsf{BVM}^{\leftarrow}(\mathcal S)\) for the
inverse-image-directed bicategory: a 1-cell from \(\mathcal E\) to
\(\mathcal F\) is a finite-limit-preserving left adjoint
\(L:\mathcal E\to\mathcal F\), with its right adjoint and a specified
base isomorphism \(Lp^*\cong q^*\). A 2-cell is a natural
transformation of these left adjoints compatible with that isomorphism.
Composition has the usual composite base isomorphism.
In the conventional geometric direction a 1-cell
\(\mathcal E\to\mathcal F\) instead has inverse image
\(\mathcal F\to\mathcal E\). We keep this opposite direction
separate. The article names a bicategory on p.~333 without printing its
1-cell or 2-cell definitions; the notation here specifies a reconstruction,
not a quotation of an unprinted convention.
\end{hbdefinition}

\begin{hbtheorem}{thm:p10-lawvere-bvm-correspondence}{proved=annex:p10.lawvere-bvm.correspondence}{Maps, adjunctions, and coherence of the correspondence}
There are biequivalences, with the algebra categories regarded as locally
discrete bicategories,
\[
 \mathsf{BVM}^{\leftarrow}(\mathcal S)
       \simeq\mathsf{CBA}(\mathcal S),\qquad
 \mathsf{BVM}_{\mathrm{geom}}(\mathcal S)
       \simeq\mathsf{CBA}(\mathcal S)^{\mathrm{op}}.
\]
A complete homomorphism \(\alpha:B\to C\) acts on Boolean matrices
entry by entry. Its right adjoint on sheaves is the actual functor
\(R_\alpha(Y)(b)=Y(\alpha b)\). Conversely a 1-cell induces its
homomorphism by transport of subobjects of base-indexed objects.
The adjunction bijections, units and counits, reconstruction isomorphisms,
uniqueness of compatible 2-cells, and composition coherence are established
in the annex. Such inverse images are geometric; they need not preserve
power objects or be logical morphisms.
\end{hbtheorem}

For comparison, Mac Lane--Moerdijk explicitly separate the internal
Boolean condition, two-valuedness, and the splitting of all epimorphisms
\cite[pp.~270--276]{MacLaneMoerdijk1992}. Their later choice-independence
construction is not a substitute for the source's stronger Boolean
package \cite[pp.~291--295]{MacLaneMoerdijk1992}; the proof of that later
construction continues beyond this selected reading. The correspondence
above is relative to the given abstract-set base, not a classification
of all elementary topoi or of ordinary set-theoretic models.
The full-datum closed functor and nondegenerate two-valued collapse
follow in Theorems~\ref{thm:p10-lawvere-closed-base-partials}--
\ref{thm:p10-lawvere-closed-base-collapse}, after this BVM construction.
The countable-chain-condition/CH transfer, global-spectrum universal
property, and independently owed recursion/free-algebra dependencies
remain separate; neither choice nor the correspondence proves them.

% Adopted Lawvere continuation; working spec fragment preserved separately.
\subsubsection{The full-datum closed-base construction and collapse}

The original p.~332 theorem distinguishes a closed-category base from a
geometric base \cite[p.~332]{Lawvere1970}; the event is Nice 1970 and the
proceedings publication is 1971. The following construction retains the
given finite-limit functor \(u_*:\mathcal E\to\mathcal E_0\) and
Heyting homomorphism \(h:u_*T\to T_0\). It does not assume that
\(u_*\) itself descends through the selected-mono localization.

For the historical terminology, a closed functor has a functor, a unit
map, and a natural internal-hom comparison satisfying identity, unit and
composition compatibilities. The hom comparison need not be invertible.
Eilenberg--Kelly distinguish this from \emph{normality}, which identifies
the underlying-element functors \cite[pp.~433--434,439]{EilenbergKelly1966}.
Their relative-base paragraph imposes a stronger normal convention
\cite[p.~549]{EilenbergKelly1966}. Lawvere's six-page article does not
print such an additional element-functor requirement. Here we construct
the actual closed functor, its base-valued hom objects and comparison;
we do not silently identify the resulting enriched global underlying
category with the ordinary localization. The countercase below explains
why that stronger identification cannot follow from the full printed datum.

\begin{hbdefinition}{def:p10-lawvere-closed-base-partials}{The h-total partial-element quotient}
Let \(\partial X\subseteq\mathcal P_{\mathcal E}X\) be the object
of subsingleton predicates, with universal partial member
\(D_X\hookrightarrow\partial X\) and evaluation \(D_X\to X\).
Write \(d_X:\partial X\to T\) for its domain classifier. On
\(\partial X\times\partial X\), let
\(e_X(p,q)\) classify the domain on which the two partial members
have a common value. Thus
\[
 e_X(p,p)=d_X(p),\quad e_X(p,q)=e_X(q,p),\quad
 e_X(p,q)\le d_X(p)\wedge d_X(q),\quad
 e_X(p,q)\wedge e_X(q,r)\le e_X(p,r).
\]
Define in \(\mathcal E_0\)
\[
 \widehat d_X=h\,u_*d_X,\qquad \widehat e_X=h\,u_*e_X,
 \qquad A_X=\{p\in u_*\partial X:\widehat d_X(p)=\mathsf{true}_0\}.
\]
The finite-product identifications of \(u_*\) are understood.
On \(A_X\) let \(R_X\) be classified by \(\widehat e_X\),
and put \(H_hX=A_X/R_X\), using the established effective quotient
construction in \(\mathcal E_0\). The hats classify newly defined
predicates, not the images under \(u_*\) of the original subobjects.
Every formula is internal in the base and applies at all generalized
parameter objects.
\end{hbdefinition}

\begin{hbtheorem}{thm:p10-lawvere-closed-base-partials}{proved=annex:p10.lawvere-closed-base.partials}{The quotient functor and its exact finite-limit behavior}
The relation \(R_X\) is an equivalence relation. Forward image of
partial members defines an actual functor
\(H_h:\mathcal E\to\mathcal E_0\), with a natural singleton map
\(\eta_X:u_*X\to H_hX\). The quotient has the universal property
that maps \(H_hX\to Z\) correspond uniquely to maps
\(A_X\to Z\) equal on \(R_X\).
There is also an actual map \(\tau:H_hT\to T_0\) with
\(\tau\eta_T=h\); it is not asserted a classifier isomorphism.

The functor preserves terminal objects, products and equalizers, with
explicit inverse product maps and equalizer lifts obtained by restricting
partial members. It inverts every mono in the original class
\(\mathcal W_h\), not merely principal or finite-filter examples.
Consequently the exact ordinary localization property yields
\[
 \overline H_h:\mathcal E_h\longrightarrow\mathcal E_0,
 \qquad
 \overline H_h[U,f]=H_h(f)H_h(i_{U,X})^{-1}.
\]
This extension preserves finite limits of arbitrary localized arrows.
It is a new base functor, not a claimed factorization of \(u_*\).
\end{hbtheorem}

\begin{hbtheorem}{thm:p10-lawvere-closed-base-closed}{proved=annex:p10.lawvere-closed-base.closed}{Actual hom objects, composition and the closed comparison}
The cartesian closed structures of \(\mathcal E_h\) and
\(\mathcal E_0\) give the unit
\(1_0\cong\overline H_h1\) and natural comparison
\[
 \kappa_{X,Y}:\overline H_h(Y^X)
                  \longrightarrow(\overline H_hY)^{\overline H_hX},
\]
obtained by transposing transported evaluation. They satisfy all three
closed-functor compatibilities. In particular
\(\mathcal C_h(X,Y)=\overline H_h(Y^X)\) are actual base hom
objects, with identities transported from names of identity maps and
composition transported from internal composition. Associativity, unit
laws and the compatibility of \(\kappa\) with composition are proved
as equations of maps in \(\mathcal E_0\).
The closed comparison is not asserted an isomorphism or normal. The
ordinary maps of \(\mathcal E_h\) have natural names in these hom
objects; the naming maps need not be bijections onto their global elements.
The quotient and localization universal properties, and the exponential
transpose property fixing this product-induced comparison, are explicit
in the annex. No adjoint or geometric direct image is assumed for
\(\overline H_h\).
\end{hbtheorem}

\begin{hblemma}{thm:p10-lawvere-closed-base-countercases}{proved=annex:p10.lawvere-closed-base.countercases}{Two different base requirements that cannot be conflated}
For product \(u_*:\mathsf{Set}\times\mathsf{Set}\to\mathsf{Set}\)
and first-coordinate truth homomorphism, the new \(H_h\) is first
projection, whereas \(u_*\) cannot descend unchanged.
For diagonal \(u_*:\mathsf{Set}\to\mathsf{Set}\times\mathsf{Set}\)
and the identity truth homomorphism, \(\mathcal E_h=\mathsf{Set}\)
and \(H_h\) is diagonal. Its base hom object is
\((Y^X,Y^X)\); its global elements are pairs of functions, while
ordinary localized maps give only equal pairs. Thus even an invertible
closed comparison does not make the enriched global underlying category
the specified ordinary localization. Neither stronger requirement is
inserted into the source theorem.
\end{hblemma}

\begin{hbtheorem}{thm:p10-lawvere-closed-base-collapse}{proved=annex:p10.lawvere-closed-base.collapse}{The nondegenerate Boolean-valued two-valued collapse}
Let \(p:\mathcal E\to\mathcal S\) be a \emph{nondegenerate}
source Boolean-valued model over the abstract-set base. Then its internally
complete Boolean algebra \(B=p_*T\) admits a Boolean homomorphism
\(h:B\to2=T_{\mathcal S}\). This need not preserve arbitrary joins.
For this original datum \((p_*,h)\), the exact localization
\(\mathcal E\to\mathcal E_h\) is logical and its target is
nondegenerate, classical, two-valued, and still satisfies choice.
It has the constructed closed functor to \(\mathcal S\); it is
not claimed geometrically defined over \(\mathcal S\). In this
particular collapse case the naming maps are bijective, giving the
normal element-functor identification up to canonical isomorphism;
this uses the geometric original \(p_*\) and the proved choice
in \(\mathcal S\), unlike the unrestricted full-datum case.

The annex constructs \(h\) by bounded well-ordered proper-filter
histories, rather than assuming an ultrafilter lemma or a ZFC universe.
Nondegeneracy is necessary: the one-element Boolean algebra admits no
Boolean homomorphism to the nondegenerate algebra \(2\). The general
full-datum construction above includes the improper-filter/degenerate-base
boundary, but that boundary is not a nondegenerate two-valued model.
\end{hbtheorem}

These are independent editorial proofs. The original closed-base and
collapse statements remain historically announcements on p.~332. They
do not settle the unprinted normal-relative convention by inventing an
element-functor identity. They also do not supply the outstanding suitable
ccc/CH comparison and transfer, spectrum/locality/universality, or the
independent recursion-lifting/free-algebra arguments.


\subsubsection{Double-negation Boolean-valued models and the CH sketch}

The article also gives the construction
\[
           \widetilde P=(P\mathcal S)_{\neg\neg}
\]
for a poset \(P\) internal to \(\mathcal S\)
\cite[p.~333]{Lawvere1970}. The BVM correspondence and collapse
just constructed do not prove the subsequent cardinal-comparison
or chain-condition claims.

The forcing notation in this construction is stagewise and explicitly
relative to \(\widetilde P\).  If \(q\geq p\), \(x\) is
an element of \(X\) defined at \(p\), and
\(\varphi:X\to2\), then \(q\) forces \(\varphi(x)\)
when \(\varphi(x/q)=\mathsf{true}\).  In double-negation
sheaves this holds at \(q\) exactly when it holds on a set of
stages cofinal beyond \(q\).  This is the source's specialization
of sheaf locality; it is not a general current-stage witness clause
for all presheaves.

\begin{hbtheorem}{thm:p10-lawvere-sheaves-ch}{sketched=annex:p10.lawvere-sheaves.ch}{Lawvere's Boolean-valued CH-refutation construction}
In \(\mathcal S\), choose \(I\) with \(2^\omega<I\), where
the inequality means a monomorphism \(2^\omega\to I\) but no
epimorphism in that same direction.  Let \(P\)
be finite partial maps \(\omega\times I\to2\), ordered by
extension, and use its double-negation sheaf topos
\(\widetilde P\).  The source calls this finite-partial-map
object intrinsically definable; its application here is over the
abstract-set base with the stated recursion object.  For the
constant-sheaf inverse image
\(u^*:\mathcal S\to\widetilde P\), the source asserts
\[
                    \omega<u^*(2^\omega)<2^\omega
                      \quad\text{in }\widetilde P.
\]
The rightmost power object is formed in \(\widetilde P\);
the middle one is the transported ground object.  The source
produces a sheaf-local injection \(u^*(I)\hookrightarrow2^\omega\)
and depends on its following lemma: for a poset satisfying a suitable
countable-chain condition and ground objects \(X,Y\) with
\(Y\times\omega\cong Y\),
\[
 \operatorname{Epi}(X,Y)=0\text{ in }\mathcal S
 \quad\Longrightarrow\quad
 \operatorname{Epi}(u^*X,u^*Y)=0\text{ in }\widetilde P.
\]
Here \(\operatorname{Epi}(X,Y)\hookrightarrow Y^X\) is the
object of maps whose image is all of \(Y\), obtained by pulling
back \(\mathsf{true}\) along the predicate of having full image.
The article leaves the word ``suitable'' in the chain-condition
hypothesis and supplies no proof of this lemma
\cite[p.~333]{Lawvere1970}. The following editorial continuation
gives a precise sufficient internal ccc and proves it for this actual
finite-partial-map object. Its no-epi, local-injection and strict
internal comparison proofs are
\cref{thm:p10-lawvere-ch-chain-ccc,thm:p10-lawvere-ch-chain-no-epi,
thm:p10-lawvere-ch-chain-generic,thm:p10-lawvere-ch-chain-comparisons},
with their cardinal and internal-site dependencies supplied separately.
This is not a claim that the article prints that chosen ccc formula
or those complete proofs. The original statement retains its historical
sketch standing. Specified set-theoretic transfer/model interpretation
and CH's positive consistency direction still require separate proofs.
\end{hbtheorem}

% Adopted CH-chain continuation; working spec fragment preserved separately.
\subsubsection{The finite-partial-map chain condition and the source's no-epi lemma}

The original Nice 1970 article, published in 1971, leaves the word
\emph{suitable} in its countable-chain-condition lemma on p.~333
\cite[p.~333]{Lawvere1970}. Here is an explicit sufficient internal
condition and its proof for the very finite-partial-map object used there.
These are independent editorial proofs, not unprinted formulas attributed
to the original source. The base is the declared abstract-set topos
\(\mathcal S\), with classical logic, choice, power objects and the source's
recursion object \(\omega\), not an external universe of ZFC sets.

\begin{hbdefinition}{def:p10-lawvere-ch-chain-poset}{Finite conditions and internally countable antichains}
Let \(D=\omega\times I\). Write \(\operatorname{Fin}(D)\) for its
finite-subset object, the least
subobject of \(\mathcal P D\) containing empty and singleton predicates
and closed under binary union, obtained by intersection of all such
subobjects. In the present classical recursion base these are exactly
subsets bijective with some initial segment \([n]\subseteq\omega\).
Let \(P\) be the object of functional relations \(D\rightharpoonup2\)
whose domain lies in this finite-subset object.

Write \(q\succeq p\) when \(q\) extends \(p\), the source's \(q\ge p\).
The site category has arrows \(q\to p\) in that case, so its presheaf
restriction runs toward stronger conditions. Conditions are compatible
iff they agree on their common domain; their union is then a common
extension. An antichain is a subobject of pairwise incompatible distinct
conditions. Countable means admitting a monomorphism into \(\omega\),
including finite and empty objects. The internal ccc says that every
antichain subobject is countable. This is an internal quantified statement
on the relevant power objects, valid with parameters; it is not merely an
external assertion about global conditions.

On any such poset use the double-negation coverage: a sieve at \(p\)
covers iff for every \(q\succeq p\) there is \(r\succeq q\) in the sieve.
Let \(\widetilde P\) be its internal sheaf topos and
\(u^*=a\Delta:\mathcal S\to\widetilde P\) its constant-sheaf inverse image.
The earlier internal-presheaf and reflector constructions supply this
topos and geometric adjunction. Constant sheaves are not assumed to have
only ground sections at every stage.
\end{hbdefinition}

\begin{hblemma}{thm:p10-lawvere-ch-chain-cardinal-tools}{proved=annex:p10.lawvere-ch-chain.cardinal-tools}{The bounded cardinal tools actually used}
In \(\mathcal S\), countable unions of internally countable families
are countable, images of countable objects are countable, and
\(\omega\times\omega\cong\omega\). Every infinite object \(X\) satisfies
\(X\times\omega\cong X\).
Pairwise-disjoint subfamilies and antichains have maximal subfamilies.
The annex constructs these by bounded well-ordered histories in power
objects and ordinary natural-number recursion, not external ordinal
recursion, an unproved Zorn lemma, or Replacement.
\end{hblemma}

\begin{hbtheorem}{thm:p10-lawvere-ch-chain-ccc}{proved=annex:p10.lawvere-ch-chain.ccc}{The actual finite-partial-map poset is internally ccc}
For every \(I\in\mathcal S\), the object \(P\) above satisfies the
internal ccc. The proof constructs an uncountable delta-system from an
uncountable family of distinct finite supports, then fixes the finite
root assignment. Equal root assignments make two conditions compatible.
Finite support fibers and countable unions are proved, not assumed to
preserve cardinal size.
\end{hbtheorem}

\begin{hblemma}{thm:p10-lawvere-ch-chain-constant-sections}{proved=annex:p10.lawvere-ch-chain.constant-sections}{Constant sections, deciding covers and number objects}
For a base object \(A\), the constant presheaf is separated and its
unit into \(u^*A\) is a dense mono. Thus every section of \(u^*A\)
is locally a ground section, and unequal ground sections never become
equal. These assertions hold in every generalized parameter context.
The image \(u^*\omega\), with transported zero and successor, is the
natural-number object of \(\widetilde P\).
This classical sheaf topos is a Boolean-valued model in the source's
stronger sense: its canonical free-action covers are subobjects of
base-indexed constant sheaves, and the proved matrix disjointification
therefore supplies choice and support splitting.
No enough-points assumption or logical preservation of arbitrary
function objects is made.
\end{hblemma}

\begin{hbtheorem}{thm:p10-lawvere-ch-chain-no-epi}{proved=annex:p10.lawvere-ch-chain.no-epi}{A precise sufficient no-epi preservation lemma}
Let an internal poset \(P\) satisfy the internal ccc just defined, and
let \(X,Y\in\mathcal S\), retaining the source's hypothesis
\(Y\times\omega\cong Y\). Then
\[
 \operatorname{Epi}_{\mathcal S}(X,Y)=0
 \quad\Longrightarrow\quad
 \operatorname{Epi}_{\widetilde P}(u^*X,u^*Y)=0.
\]
Here the epi objects are the full internally parametrized predicates of
surjectivity, not the absence of one chosen global map. For each ground
input, an actual maximal deciding antichain gives a countable family of
possible output values. A purported local epi yields a ground epi from
\(X\times\omega\); infinite \(X\) absorbs that factor. Finite \(X\)
is handled by deciding all its values on a common extension and choosing
a missing ground output. The proof does not make the invalid inference
that a surjection from \(X\times\omega\) always gives one from \(X\).
\end{hbtheorem}

\begin{hbexample}{def:p10-lawvere-ch-chain-naive-countercase}{Why the finite-input branch is necessary}
For \(X=1,Y=\omega\), the projection \(X\times\omega\to Y\) is
an isomorphism, while no map \(X\to Y\) is epi: a singleton image
cannot contain both zero and successor of zero. Also
\(Y\times\omega\cong Y\). Thus even the source's absorption
hypothesis does not justify the naive intermediate cardinal inference.
The proved finite-input deciding argument is what repairs that step.
\end{hbexample}

\begin{hbtheorem}{thm:p10-lawvere-ch-chain-generic}{proved=annex:p10.lawvere-ch-chain.generic}{Actual generic evaluation and local injection descent}
For the finite-partial-map \(P\), the domain where a condition assigns a
value at \((n,i)\) is a dense subpresheaf of \(\Delta(\omega\times I)\).
Its evaluation into \(\Delta2\), after sheafification, gives an actual map
\[
 u^*I\times u^*\omega\longrightarrow u^*2.
\]
Its transpose \(\alpha:u^*I\to2^\omega\), using the new sheaf topos's
power object and number object, is monic. The proof uses a fresh finite
coordinate to distinguish two unequal indices and proves the required
density and sheaf descent at arbitrary parameter bases.
\end{hbtheorem}

\begin{hbtheorem}{thm:p10-lawvere-ch-chain-comparisons}{proved=annex:p10.lawvere-ch-chain.comparisons}{The displayed strict comparisons at their internal scope}
Put \(C=2^\omega\) in \(\mathcal S\), and take \(I\) with a mono
\(C\to I\) and \(\operatorname{Epi}_{\mathcal S}(C,I)=0\), as in the
source. Such an \(I\) is available, for example \(I=2^C\), by the
singleton injection and Cantor's diagonal argument.
In the finite-partial-map sheaf topos the actual comparisons are
\[
             \omega<u^*C<2^\omega.
\]
Each strict inequality means an exhibited mono in that direction and
zero internally parametrized epi object in the same direction.
The middle-to-right mono is the transported-evaluation comparison, not
an assumed isomorphism \(u^*(2^\omega)\cong2^\omega\).
The zero epi object on the right follows from the generic injection and a complement
retraction, together with the proved no-epi lemma; no extra global sheaf
choice is assumed for that retraction.

This supplies the source's internal cardinal comparison and its omitted
ccc/no-epi/local-injection dependencies. It does not specify or prove
a translation to a separately chosen set-theoretic axiom system or model,
nor CH's positive consistency direction. Spectrum and the independently
owed general recursion/free-algebra constructions remain separate.
\end{hbtheorem}


\subsubsection{The ringed-topos global spectrum}

The third source section first takes a commutative multiplicative
monoid object \(A\) and a homomorphism
\(u:A\to T^X\), where multiplication in \(T^X\) is
intersection.  Its pairing defines a membership subobject
\(P\hookrightarrow X\times A\), and the inclusion relation
between basis elements defines \(F\rightrightarrows A\).
The latter gives a submonoid of the constant monoid presheaf
\(\overline A\) on the index category; the former gives a
submonoid of the constant family \(\dot A\) over \(X\).
Forming fractions gives \((\overline A)_F\) on the presheaf
side and \((\dot A)_P\) on the family side.  The intermediate
sheaf category has the reflected monoid \(\widetilde A\)
between them \cite[p.~333]{Lawvere1970}.

\begin{hbdefinition}{def:p10-lawvere-sheaves-prime}{The source's prime subobject and local ring}
For a commutative ring object \(A\), a prime \(x\) in the
source's terminology is a subobject of \(A\), not an ideal,
whose membership predicate satisfies, internally,
\[
\begin{aligned}
 [1\in x]&=\mathsf{true},\\
 [fg\in x]&=[f\in x]\land[g\in x],\\
 [0\in x]&=\mathsf{false},\\
 [f+g\in x]&\leq[f\in x]\lor[g\in x].
\end{aligned}
\]
The multiplication clause is an equivalence, not one implication.
The disjunction is the join of subobjects, and may hold only
locally.  The source notes that this intuitionistic distinction
already matters in \(2\mathsf{Set}\).  A ring is local when
its subobject of units is prime in
this sense.  The object \(X\hookrightarrow T^A\) of such
subobjects is obtained by imposing these conditions with finite
limits on predicate objects and their quantified maps
\cite[pp.~333--334]{Lawvere1970}.
\index{prime subobject!Lawvere's spectrum}
\index{local ring!internal definition}
\end{hbdefinition}

\begin{hbtheorem}{thm:p10-lawvere-sheaves-spectrum}{sketched=annex:p10.lawvere-sheaves.spectrum}{The announced global-spectrum universal construction}
The resulting pairing gives a topology basis whose sheaves form
\(\operatorname{Spec}(A)\), the global spectrum of the ringed
topos \((\mathcal E,A)\).  In it \(\widetilde A\) is a local
ring object and is the universal local \(A\)-algebra among
topoi defined over \(\mathcal E\).  The article describes the
membership relation as being transformed into its opposite in this
process; the fractions invert the selected multiplicative subobjects.
This is the source's intrinsic
Chevalley--Hakim construction \cite[p.~334]{Lawvere1970}.
The article outlines the monoid, fraction, and sheaf steps; it does
not supply their universal-property proofs.  The annex keeps those
dependencies explicit.
\end{hbtheorem}

\begin{hbexample}{def:p10-lawvere-sheaves-classical-prime}{Comparison with the complement of a prime ideal}
In ordinary sets with classical logic, the four conditions say that
\(x\) is the complement of a proper prime ideal.  For
\(A=\mathbb Z\) and a prime number \(\ell\), take
\(x=\{n:\ell\nmid n\}\).  Its localization inverts just
these elements and gives \(\mathbb Z_{(\ell)}\); the units
are its localized image of \(x\).  This worked comparison is
editorial.  In intuitionistic internal logic, replacing the source's
prime subobject by its putative ideal complement would erase the
local-disjunction qualification.
\end{hbexample}

The closing section describes the algebraic-geometric application as
only beginning.  It attributes an interpretation of intuitionistic
semantics in a suitable varying-set topos to unpublished work of
George Rousseau, and the analogous interpretation of L\"auchli's
proof-theoretic semantics to Anders Kock.  It proposes parameters for
actual time rather than only imagined construction stages.  It says
that each definition of the reals determines an object in a topos
satisfying \((\omega)\), but records the analytic theorems then
available as unresolved.  These are historical attributions and the
article's 1970 questions, not assertions of a current unresolved
state \cite[p.~334]{Lawvere1970}.

The source bibliography identifies Mao's 1966 philosophical work;
Grothendieck--Giraud--Verdier's \emph{Topologie et faisceaux}; Hakim's
1967 Orsay thesis \emph{Sch\'emas relatifs}; Solovay--Scott's
Boolean-valued-model work as forthcoming; and Lawvere's 1969
\emph{Adjointness in Foundations} and 1970 work on equality and
comprehension in hyperdoctrines.  Those references locate the
article's dependencies; they are not substituted here for acquired
full proofs of the announced claims.

\begin{proofboundary}{def:p10-lawvere-sheaves-proof-boundary}{Source standing and the independent proof obligation}
The original article's complete account is distinct from a complete
modern development of its announced results.  The concrete and
internal presheaf packages, dependent products, internal
Heyting/predicate adjoints, stable images and effective quotients,
both internal-functor adjoints, choice equivalence, closure laws,
subterminal-generator argument, all four associated-sheaf
observations, sheaf characterization, left-exact reflector,
sheaf-topos/classifier package, geometric factorization,
double-negation classicality, and the intrinsic Godement
realization/image/cotriple reconstruction have
full, separately attributed editorial annex proofs. The exact ordinary
localization and logical preservation, Boolean-matrix classification
and choice at the declared cell regimes, and full-datum closed functor
with nondegenerate two-valued collapse likewise have their actual
editorial constructions above. General strong or normal closed-base
identification is not inferred from the arbitrary original datum;
the diagonal countercase and the separate normality proof for the
collapse case retain that distinction.

The explicit sufficient internal ccc, its finite-partial-map verification,
cardinal tools, internal-site bridge, complete no-epi predicate argument,
local injection descent and both strict internal CH comparisons now have
the six reciprocal editorial proofs in the continuation above. The
unprinted word suitable is not retroactively assigned a unique original
formal definition, and the proofs retain all declared base parameters
and cardinal/choice/NNO hypotheses.

Specified set-theoretic transfer/model interpretation, CH's positive
consistency direction, spectrum locality/universality and the independent
general recursion-lifting/free-algebra dependencies remain owed. Arbitrary
strong or normal-base obligations are not discharged by this internal
forcing argument. Historical source announcement or sketch does not
waive a separately promised full proof. The reciprocal source sketch
targets retain these actual remaining and historical boundaries.
\end{proofboundary}
